
\medskip
\noindent\textbf{4.2.}
\phantomsection
\label{sga3:VIIA:4.2}

We now consider a construction closely related to the preceding
one. Let \(S\) be a scheme of characteristic \(p\), let \(X\) be
an \(S\)-scheme, and let \(X_S^p\) be the product, in
\((\mathbf{Sch}/S)\), of \(p\) copies of \(X\).

Denote by \(U^p(X)\) the open subscheme of \(X_S^p\) which is the
union of the products \(U_S^p\) as \(U\) ranges over the affine
open subsets of \(X\). Thus a point \(x\) of \(X_S^p\) belongs to
\(U^p(X)\) if and only if its projections
\(\operatorname{pr}_i x\) on the factors of \(X_S^p\) belong to a
common affine open subset of \(X\). For example, if every finite
subset of \(X\) is contained in an affine open subset, then
\[
  U^p(X)=X_S^p.
\]

The symmetric group \(\mathfrak S_p\) on \(p\) letters acts on
\(X_S^p\) by permuting the factors and leaves \(U^p(X)\) stable.
We call the quotient of \(X_S^p\) by \(\mathfrak S_p\) in the
category of ringed spaces the \emph{\(p\)-fold symmetric product
of \(X\)}, and denote it by \(\Sigma^p X\). Let
\[
  q(X),\ \text{or simply }q:
  X_S^p\longrightarrow\Sigma^p X
\]
be the canonical projection.

The morphism \(q\) maps \(U^p(X)\) onto an open subset
\(V^p(X)\) of the symmetric product, which may be described as
follows (cf.\ V, \SGARefLink{thm:sga3-V-4.1}{4.1}). The structure sheaf of \(\Sigma^p X\)
induces a scheme structure on \(V^p(X)\), and the morphism
\[
  q'(X):U^p(X)\longrightarrow V^p(X)
\]
induced by \(q(X)\) is affine, indeed integral. As \(U\) ranges
over the affine open subsets of \(X\) which map into a varying
affine open \(V\) of \(S\), the schemes \(\Sigma^p U\) form an
affine open covering of \(V^p(X)\). If \(R\) is the affine algebra
of \(V\), and \(A\) that of \(U\), the affine algebra of
\(\Sigma^p U\) is the subalgebra \(\Sigma^p A\) of
\(\bigotimes_R^p A\) consisting of the symmetric tensors.

Now consider the diagonal morphism
\[
  \delta:X\longrightarrow U^p(X).
\]
If \(V=\operatorname{Spec}R\) is an affine open subset of \(S\)
and \(U=\operatorname{Spec}A\) is an affine open subset of \(X\)
lying over \(V\), the restriction of \(\delta\) to \(U\) is
defined by the algebra homomorphism
\[
  \eta:\bigotimes_R^p A\longrightarrow A,
  \qquad
  a_1\otimes\cdots\otimes a_p
  \longmapsto
  a_1a_2\cdots a_p.
\]
Consequently, if \(N\) denotes the symmetrization operator,
\[
\begin{aligned}
  \eta\bigl(N(a_1\otimes\cdots\otimes a_p)\bigr)
  &=
  \eta\left(
    \sum_{\sigma\in\mathfrak S_p}
      a_{\sigma(1)}\otimes\cdots\otimes a_{\sigma(p)}
  \right) \\
  &=
  p!\,a_1\cdots a_p
  =
  0.
\end{aligned}
\]
In other words, \(\eta\) vanishes on the subspace
\[
  N\!\left(\bigotimes_R^p A\right)
  \subseteq
  \Sigma^p A
\]
consisting of symmetrized tensors. Moreover, if \(f\) is a
symmetric tensor, then plainly \(N(fa)=fN(a)\). Hence
\(N(\bigotimes_R^p A)\) is an ideal of \(\Sigma^p A\). We put
\[
  U^{[p/S]}
  =
  \operatorname{Spec}\left(
    \Sigma^p A
    \big/
    N\!\left(\bigotimes_R^p A\right)
  \right).
\]
This is a closed subscheme of
\(\Sigma^p(U)=V^p(U)\). As \(U\) ranges over the affine open
subsets of \(X\) which map into a varying affine open \(V\) of
\(S\), the schemes \(U^{[p/S]}\) glue to a closed subscheme of
\(V^p(X)\), denoted \(X^{[p/S]}\).

If \(i(X)\) denotes the inclusion of \(X^{[p/S]}\) in \(V^p(X)\),
the preceding calculation shows that
\(q'(X)\circ\delta\) factors through \(X^{[p/S]}\). We therefore
obtain a morphism
\[
  F^{[p]}(X/S):X\longrightarrow X^{[p/S]}.
\]
The factorization is displayed by the commutative diagram
\begin{center}
  \phantomsection
  \label{sga3:VIIA:4.2-pfold-symmetric-factorization}
  \begin{tikzcd}[row sep=3.2em,column sep=1.8em]
X_S^p \arrow[dd,swap,"q(X)"]
  & \supset
  & \mathrm U^p(X) \arrow[dd,swap,"q'(X)"]
  && X \arrow[ll,swap,"\delta(X)"] \arrow[dd,"F^{[p]}(X/S)"] \\
\\[-1.2em]
\Sigma^p(X)
  & \supset
  & \mathrm V^p(X)
  && X^{[p/S]} . \arrow[ll,swap,"i(X)"]
\end{tikzcd}

\end{center}
It is clear that \(X^{[p/S]}\) is functorial in \(X\), and that
\[
  F^{[p]}:X\longmapsto F^{[p]}(X/S)
\]
is a functorial homomorphism.

\medskip
\noindent\textbf{4.2.1.}
\phantomsection
\label{sga3:VIIA:4.2.1}

The schemes \(X^{[p/S]}\) and \(X^{(p/S)}\) are naturally related.
Let \(V\) be an affine open subset of \(S\), with affine algebra
\(R\), and let \(U\) be an affine open subset of \(X\) lying over
\(V\), with affine algebra \(A\). If \(\pi\) denotes the
endomorphism \(x\mapsto x^p\) of \(R\), the affine algebra of
\(U^{(p/S)}\) is \(A\otimes_\pi R\). The map
\[
  a\otimes_\pi\lambda
  \longmapsto
  \left(
    \lambda a\otimes\cdots\otimes a
    \bmod
    N\!\left(\bigotimes_R^p A\right)
  \right)
\]
is a homomorphism of \(R\)-algebras
\[
  A\otimes_\pi R
  \longrightarrow
  \Sigma^p A\big/
    N\!\left(\bigotimes_R^p A\right).
\]
It induces a morphism
\[
  \varphi(U):U^{[p/S]}\longrightarrow U^{(p/S)}
\]
satisfying
\[
  \varphi(U)\circ F^{[p]}(U/S)=\operatorname{Fr}(U/S).
\]
Gluing these morphisms gives the commutative triangle
\begin{center}
  \phantomsection
  \label{sga3:VIIA:4.2.1-phi-frobenius-triangle}
  \begin{tikzcd}[row sep=3.6em,column sep=4.4em]
& X
  \arrow[dl,swap,"F^{[p]}(X/S)"]
  \arrow[dr,"\operatorname{Fr}(X/S)"] & \\
X^{[p/S]} \arrow[rr,"\varphi(X)"]
  && X^{(p/S)} .
\end{tikzcd}

\end{center}

For example, if \(X\) is the subscheme of \(S\) defined by a
quasi-coherent ideal \(\mathcal I\), then
\(F^{[p]}(X/S)\) identifies with the identity morphism of \(X\).
Thus \(\varphi(X)\) is the canonical immersion
\[
  \operatorname{Spec}(\mathcal O_S/\mathcal I)
  \longrightarrow
  \operatorname{Spec}(\mathcal O_S/\mathcal I^{\{p\}}).
\]
This shows that \(\varphi(X)\) need not be an isomorphism.

On the other hand, when \(M\) is a free \(R\)-module, the map
\[
\begin{aligned}
  M\otimes_\pi R
  &\longrightarrow
  \Sigma^p M\big/
    N\!\left(\bigotimes_R^p M\right), \\
  m\otimes_\pi\lambda
  &\longmapsto
  \left(
    \lambda m\otimes\cdots\otimes m
    \bmod
    N\!\left(\bigotimes_R^p M\right)
  \right)
\end{aligned}
\]
is plainly bijective. It remains bijective when \(M\) is
\(R\)-flat, because every flat module is a filtered direct limit of
free modules.
It follows that
\[
  \boxed{\;
  \varphi(X):X^{[p/S]}\longrightarrow X^{(p/S)}
  \text{ is an isomorphism if \(X\) is flat over \(S\).}
  \;}
\]

